\chapter{Umsetzung}\label{chap:approach}
\section{Komponenten des Verfahrens}
Das vorgestellte Verfahren besteht aus den folgenden Komponenten:
\begin{itemize}
    \item testDiscovery: liefert die Gesamtheit aller Testcases
    \item output: Enthält die Ausgabe aller bisher durchführten Testdurchläufe
    \item \textit{testAuditor}: Evaluiert, welche Testcases noch ausgeführt werden müssen
    \item testExecution: Führt eine gegebene Menge von Testcases aus, befüllt output
    \item manager: Steuert die Interaktion mit dem user, sowie die Kommunikation von  Komponenten untereinander
    \item user: Interagiert nach eigenem Ermessen mit dem manager
\end{itemize}
Eine Übersicht zur Interaktion der einzelnen Komponenten liefert \figref{fig:tool-overview}:
\begin{figure}[htbp]
    \centering
    \begin{tikzpicture}[
        node distance=1.1cm and 2.0cm,
        box/.style={
            draw,
            rectangle,
            thick,
            minimum width=3.4cm,
            minimum height=1.1cm,
            align=center
        },
        smallbox/.style={
            draw,
            rectangle,
            thick,
            minimum width=3.1cm,
            minimum height=0.9cm,
            align=center
        },
        botharrow/.style={
            <->,
            >=Latex,
            thick
        }
    ]

        % Zentrale Komponente
        \node[box] (script) {Script};
        % Linke Eingaben
        % above left=1.0cm and 1.0cm
        \node[smallbox, left=1.2cm of script] (ci) {CI};
        \node[smallbox, below left=1.2cm and 1cmof script] (dev) {Entwickler\\(lokal)};

        % Rechte Komponenten
        \node[smallbox, right=2.2cm of script, yshift=1.4cm] (notes) {git notes/\\sonstiges};
        \node[smallbox, right=2.2cm of script, yshift=-0.4cm] (discovery) {testDiscovery.sh};
        \node[smallbox, below=0.5cm of discovery] (auditor) {testAuditor};
        \node[smallbox, below=0.5cm of auditor] (execution) {testExecution.sh};

        % Bidirektionale Verbindungen links
        \draw[botharrow] (ci.south) -- ++(0,-0.45) -| (script.north west);
        \draw[botharrow] (dev.east) -- (script.west);

        % Bidirektionale Verbindungen rechts
        \draw[botharrow] (script.north east) -- (notes.west);
        \draw[botharrow] (script.east) -- (discovery.west);
        \draw[botharrow] (script.south east) -- (auditor.west);
        \draw[botharrow] (script.south) -- (execution.west);
    \end{tikzpicture}
    \caption{Schematische Darstellung der Interaktion zwischen CI, lokalem Entwickler, Script und Hilfskomponenten}
    \label{fig:tool-overview}
\end{figure}

Im Folgenden werden die einzelnen Komponenten des Verfahrens, so wie ihr Zusammenspiel untereinander im Detail erläutert.

\subsection{testDiscovery}
Durch die Wahl für JUnit XML als Ausgabeformat, das vom Verfahren verstanden werden soll, müssen die in \chapref{Evaluation_Junit} aufgelisteten Anforderungen, die von JUnit XML nicht erfüllt werden, anderweitig umgesetzt werden.\\
Die größte Herausforderung ist, dass aus einer JUnit-XML Ausgabe nur jene Testcases ermittelt werden können, welche während des Testlaufes ausgeführt wurden. \textit{testAuditor} benötigt jedoch zwingend Informationen über die Gesamtheit aller Testcases im Projekt.

Diesem Zweck dient die \textit{testDiscovery:}\\
Bei dieser Komponente ist die Implementierung als shellscript am geeignetsten, da sie bei einem Aufruf die Gesamtheit aller Testcases des Projektes in einem definierten Format zurückgibt. Manche Testframeworks bieten Zugriff auf ihre interne Discovery für Testcases, was in diesem Fällen die Implementation recht einfach gestaltet. 

\subsubsection{Format der Ausgabe der testDiscovery}
\textit{testAuditor} erwartet die Ausgabe der testDiscovery in diesem Format:
\begin{lstlisting}[caption={Format der testDiscovery}, label={TAPmin}]
<?xml version="1.0" encoding="utf-8"?>
<testsuite>
    <testcase classname="" name="" qualifiedName=""/>
    <testcase classname="" name="" qualifiedName=""/>
</testsuite>
\end{lstlisting}
Der Inhalt der beiden tags ''classname'' und ''name'' für einen Testcase muss dabei zwingend identisch mit dem Inhalt der beiden tags sein, wie er in der JUnit-XML Ausgabe des Testcase steht, da in einem späterem Schritt über die Kombination dieser beiden tags die Testcases eindeutig im JUnit XML identifiziert werden.\\
''qualifiedName'' enthält eine Zeichenfolge, über die das Testframework einen Testcase identifizieren, und ausführen kann. Je nach Testframework werden dazu der Name des Testcase, Dateipfade, Namen von klassen oder Testsuites und/oder weitere Informationen benötigt. ''qualifiedName'' dient in dem Schritt der testExecution (siehe \chapref{Komponente_testExecution}) bei der Ausführung einzelner Testcases.

Die testDiscovery ist hochgradig vom verwendeten Testframework abhängig, da die spezielle Implementierung der Testausführung des jeweiligen Testframeworks beachtet werden muss. Daher muss diese Komponente für jedes Softwareprojekt eigens angepasst werden.


\subsection{\textit{testAuditor}}\label{Komponente_testAuditor}

\subsubsection{Parser für JUnit XML}

\subsection{testExecution}\label{Komponente_testExecution}
\subsection{manager}
\subsection{user}

% The approach starts with the problem definition and continues with what you have done. Try to give an intuition first and describe everything with words and then be more formal like `Let $g$ be ...'.

% Strive to make the following items crystal clear to your readers.
% \begin{itemize}
% \item What is the problem you are treating?
% \item What are the research questions you want to answer in this work?
% \item What are your own contributions and what is work of others? If
%   you are working in a team, it is important to demarcate against the
%   contributions of others in your team.
% \end{itemize}

% Describe each major technical problem and how it was solved.

\section{Anforderungen an ein Projekt}
\subsection{Deterministische builds - 'it works on my machine'}
Damit in dem hier vorstegtellten Verfahren eine qualifizierte Aussage über die Testergebnisse treffen kann, muss gewährleistet sein, dass alle Tests unabhängig von der Entwicklungsumgebung immer das selbe Testergebnis liefern. Ein Testfall, der auf der Machine eines Entwicklers ein bestimmtes Ergebnis liefert, muss auch zwingend auf einem Build-Server das selbe Ergebnis liefern. Andernfalls könnte eine Diskrepanz zwischen den Testergebnissen zu einer verfälschten Aussage des Verfahrens über den aktuellen Stand der Test führen. 

Da Junit XML nicht alle in \chapref{Anforderungen Testausgabeformat} aufgeführten Anforderungen erfüllt, ergeben sich einige weitere Anforderungen an ein Projekt, in welchem das Verfahren angewerdet werden soll, damit alle in \chapref{Anforderungen Tool} gesetzten Ziele erreicht werden können.

\subsection{Benötigte Schnittstellen}

\subsubsection{test discovery}\label{test discovery}
Damit das Tool die Gesamtmenge aller Tests im Projekt kennt, muss im Projekt ein Script vorhanden sein, welches die die Gesamtmenge aller Tests in einem spezifiziertem Format augibt. 

-> qualifiedname: Eine Zeichenfolge, die einen Test innerhalb eines Testframeworks eindeutig identifiziert. Bespiele in vers. TestFrameworks bringen!

-> Hier Format beschreiben/definieren

\subsubsection{test execution}\label{test execution}
Zum Anderen muss dem Tool ein zweites Script zur Verfügung gestellt werden, welches eine Liste von zuvor in \chapref{test discovery} beschriebenen qualifiedNames entgegennimmt, nur diese Tests ausführt, und das Ergebnis (auch wieder JUnit XML) an das Tool zurückmeldet, welches dann diesen Input zusätzlich an die gitnotes anhängt.


\section{Anforderungen an ein Framework}
- kann JUnit XML

- enthält essenzielle Informationen für einen Test (genau ausführen!)

- kann tests einzeln aufrufen 

\section{Einschränkungen}
- Annahme: keine deliberate bad actors; keine Authentifizierung

- Annhame: An gitnotes angehängte Dateien spiegeln aktuellen Stand vom Code wieder, keine Veränderung der Codebase im Nachhinein

\section{Codestruktur}

\subsection{Parser für JUnit XML}
- Parser des tools
- leerzeilen als standart-separator von gitnotes


\section{Exemplarischer Ablauf}
% git add -A
% git commit -m "foo bar"
git notes add -F "\$(realpath code/out/report.xml)" HEAD
% git push


\section{git notes}
müssen nicht zwingend verwendet werden, auch andere Kommunikationswege mit Tool könnten etabliert werden. Hier nur Mittel der Wahl, weil ändern Code nicht.

Nicht zu empfehlen: Dateien im repo lagern.
